% Standalone insertion fragment: exact operator-theoretic consequences and scope.

\section{The operator carrier: normal form, determinants, and scope}

We first isolate the universal operator-theoretic content of the construction.
Throughout this section, inner products are linear in the second variable.

\begin{proposition}[Rank-one similarity collapse]\label{prop:rank-one-collapse}
Let $E$ be a complex Hilbert space, let $0\ne\psi,h\in E$, and put
\[
 C=|\psi\rangle\langle h|,
 \qquad \alpha=\langle h,\psi\rangle\ne0.
\]
Let $S$ be the unilateral shift on $\ell^2(\mathbb N_0)$ and define
\[
 \mathsf A=S^*\otimes C,
 \qquad
 \mathsf B=S\otimes C.
\]
There is an $R\in\operatorname{GL}(E)$ and an orthogonal rank-one
projection $P_1$ such that
\[
 RCR^{-1}=\alpha P_1.
\]
Consequently the pair $(\mathsf A,\mathsf B)$ is simultaneously similar to
\[
 ((\alpha S^*)\oplus0,(\alpha S)\oplus0).
\]
Moreover, if $P_0=|e_0\rangle\langle e_0|$, then
\[
 [\mathsf A,\mathsf B]=P_0\otimes C^2,
 \qquad
 \operatorname{Tr}[\mathsf A,\mathsf B]=\alpha^2,
\]
and hence
\begin{equation}\label{eq:universal-hhp}
 \det_F\!\left(
 e^{\mathsf A}e^{\mathsf B}e^{-\mathsf A}e^{-\mathsf B}
 \right)=e^{\alpha^2}.
\end{equation}
\end{proposition}

\begin{proof}
The identity $C^2=\alpha C$ shows that $P=C/\alpha$ is a bounded
rank-one idempotent.  The decomposition
$E=\operatorname{ran}P\dotplus\ker P$ is topological, so a bounded
invertible change of coordinates sends $P$ to an orthogonal rank-one
projection.  Tensoring this change of coordinates with the identity on
$\ell^2$ gives the asserted simultaneous normal form.  Finally,
$S^*S=I$, $SS^*=I-P_0$, and
$\operatorname{tr}(C^2)=\alpha^2$.  The last assertion is the
Helton--Howe--Pincus formula; see, for example, \cite{Migler2017}.
\end{proof}

For the Euler specialization, let
\[
 H=L^2(\mathbb R_+,e^{-x}\,dx),\qquad Xf(x)=xf(x),\qquad A=X+I,
\]
take $E=H\oplus\mathbb C$, and put
\[
 \psi=(-A^{-1}1,1),\qquad h_\gamma=(A\log X\,1,0).
\]
Both vectors belong to $E$, and direct integration gives
\[
 \langle h_\gamma,\psi\rangle
 =-\int_0^\infty e^{-x}\log x\,dx=\gamma.
\]
Thus $C=|\psi\rangle\langle h_\gamma|$ has $\alpha=\gamma$, and
\eqref{eq:universal-hhp} becomes $e^{\gamma^2}$.  The proposition also
makes the limitation precise: the simultaneous similarity class sees only
the scalar $\gamma$ and is blind to the zero-mode norm $2-\delta$ and to
the remaining Hilbert-complex geometry.

There is, however, a natural relative Fredholm determinant on the
$\delta$-side.  With the preceding $H$ and $A$, put $u=A^{-1/2}1$ and let
\[
 Q:\mathcal D(A^{1/2})\oplus\mathbb C\longrightarrow H,
 \qquad Q(f,c)=A^{1/2}f+cu.
\]

\begin{proposition}[Zero-mode perturbation determinant]
The operator $Q$ is a closed surjective Fredholm operator of index one,
\[
 QQ^*=A+|u\rangle\langle u|,
 \qquad
 \ker Q=\mathbb C\psi,
 \qquad
 \psi=(-A^{-1}1,1).
\]
For $z\in\mathbb C\setminus[1,\infty)$ define
\[
 \Delta_Q(z)=
 \det_F\!\left((QQ^*-z)(A-z)^{-1}\right).
\]
Here the product denotes its bounded extension from the natural common
domain.
Then
\begin{align}
 \Delta_Q(z)
 &=1+\langle u,(A-z)^{-1}u\rangle \notag\\
 &=1+\int_0^\infty
       \frac{e^{-x}}{(1+x)(1+x-z)}\,dx.\label{eq:relative-determinant}
\end{align}
Writing $\delta=eE_1(1)$ and taking the principal branch of $E_1$, this is
\[
 \Delta_Q(z)=1+
 \frac{e^{1-z}E_1(1-z)-\delta}{z}
 \qquad(z\ne0),
\]
with a removable value at zero given by
\begin{equation}\label{eq:delta-relative-determinant}
 \Delta_Q(0)=2-\delta=\|\psi\|^2.
\end{equation}
\end{proposition}

\begin{proof}
Surjectivity follows by taking $(f,c)=(A^{-1/2}g,0)$ for $g\in H$;
the kernel formula is immediate.  Direct computation gives
$QQ^*=A+|u\rangle\langle u|$.  The rank-one Fredholm determinant lemma
then yields \eqref{eq:relative-determinant}.  Finally,
\[
 \int_0^\infty\frac{e^{-x}}{(1+x)^2}\,dx=1-\delta,
\]
which proves \eqref{eq:delta-relative-determinant} and the index assertion.
\end{proof}

Let $\mathscr H=\ell^2(\mathbb N_0)\otimes E$, let
$\pi:\mathcal B(\mathscr H)\to
\mathcal B(\mathscr H)/\mathcal L^1(\mathscr H)$ be the quotient map, and
put $a=\pi(e^{\mathsf A})$, $b=\pi(e^{\mathsf B})$.  These elements commute.
If $\xi=\{a,b\}\in
K_2^{\mathrm{alg}}(\mathcal B(\mathscr H)/\mathcal L^1(\mathscr H))$, then
Brown's determinant invariant gives~\cite{Migler2017,Kaad2011}
\begin{equation}\label{eq:k2-boundary}
 \det\bigl(\partial_2\xi\bigr)
 =\det_F(e^{\mathsf A}e^{\mathsf B}
          e^{-\mathsf A}e^{-\mathsf B})
 =e^{\alpha^2}.
\end{equation}
In particular, when $e^{\alpha^2}\ne1$, both
$\partial_2\xi\in
K_1(\mathcal B(\mathscr H),\mathcal L^1(\mathscr H))$ and $\xi$ are
nontrivial.  Notice that \eqref{eq:k2-boundary} is an algebraic
$K_2$ statement detected on a relative $K_1$ boundary; it is not a
``relative $K_2$'' class.

\paragraph{Scope.}
The preceding statements do not produce an arithmetic relation between
$e^{\gamma^2}$ and $\delta$.  Indeed, replacing $h$ by any covector with
$\langle h,\psi\rangle=\alpha$ realizes the arbitrary scalar
$e^{\alpha^2}$, while Proposition~\ref{prop:rank-one-collapse} erases
$\delta$ completely.  Nor does \eqref{eq:k2-boundary} yield a new
topological $C^*$-algebra $K$-class: after passage to the Calkin algebra,
the two displayed invertibles are exponentials and hence have trivial
individual topological $K_1$ classes.

Under the Hardy-space identification, the active normal form is the
one-mode Toeplitz pair with symbols $e^{\alpha z^{-1}}$ and $e^{\alpha z}$.
For $\alpha\ne0$, the first symbol has an essential singularity at the
origin, so meromorphic tame-symbol formulas do not apply.  Finally, no
Wodzicki-residue interpretation follows from this model: no classical
pseudodifferential calculus or spectral triple has been specified, the
relevant commutator is finite rank and hence is annihilated by singular
traces, and a multiplicative residue determinant sends a multiplicative
commutator to $1$ whenever all factors lie in its multiplicative domain.
Thus \eqref{eq:universal-hhp} is an ordinary Fredholm
determinant anomaly, not a noncommutative-residue anomaly.
